\section{Generator and Provenance}
\label{sec:generator}

\tesserahash{} ships two API shapes: a \emph{generator} that maintains a counter for provenance-marking the same content multiple times, and a \emph{stateless} hash-only function for callers who do not need the 9-byte header.

\subsection{The stateful generator}

\begin{lstlisting}[language=Pseudocode]
g = tessera_hash.Generator()
id1 = g.new(content)             # 73 bytes: TS_now || 0x03 || ctr || SHAKE256(content)
id2 = g.new(content)             # same content, fresh TS+counter
assert id1[9:] == id2[9:]        # tails identical (same content)
assert id1 != id2                # headers differ (different TS or counter)
\end{lstlisting}

The generator is \textbf{internally locked where the language has shared-memory threading}: the \texttt{(last\_ts, counter)} read-modify-write is atomic across threads (Python \texttt{threading.Lock}, Go mutex, Rust \texttt{Mutex}, Java \texttt{synchronized}). The C struct is the documented exception: it is \emph{caller-synchronized} by contract --- \texttt{tessera\_hash.h} states the struct ``exposes no synchronization'' and callers MUST declare one generator per thread or wrap it with a mutex; the C++ wrapper inherits that contract. This matches TesseraID's \texttt{HLCGenerator} discipline (TesseraID's \texttt{StandardGenerator} is deliberately \emph{not} locked; \tesserahash{}'s generator is locked wherever the language can lock it, because counter monotonicity is the only provenance property the format offers and racing it produces undefined ordering).

The counter advances monotonically within a millisecond and resets to 0 on a new millisecond. On counter exhaustion (more than $2^{16} = 65{,}536$ identifiers generated by a single generator within one millisecond) the generator MUST spin until the wall clock strictly exceeds the last emitted timestamp (\code{now > last\_ts}), take that clock value as the new timestamp, and reset the counter to 0. If the wall clock reads earlier than \code{last\_ts} (a backward clock jump), the generator clamps to \code{last\_ts} so emitted timestamps never regress; the spin condition \code{now <= last\_ts} therefore also covers regression --- the loop cannot exit while the clock still reads at or before the last emitted timestamp. This is the same spin-to-next-millisecond overflow discipline used by TesseraID's Standard generator (which spins on exhaustion; the ``bump the logical clock by 1'' rule belongs to TesseraID's \emph{HLC} variant, not to this format).

\subsection{The stateless hash-only function}

\begin{lstlisting}[language=Pseudocode]
tail = tessera_hash.hash_only(content)         # 64 bytes: SHAKE256(content, 64)
\end{lstlisting}

\texttt{hash\_only(content)} returns just the 64-byte tail. There is no header, no timestamp, no counter. Two callers in different processes who hash the same content with this function get the same 64 bytes. Use this when:
\begin{itemize}
\item You want to deduplicate by content without committing to a generator identity.
\item You want to use the tail as an AEAD AD (e.g., binding a \tesseraseal{} envelope to specific content; see \cref{sec:monorepo}).
\item You are computing the tail for verification: given an existing \tesserahash{}, check that the trailing 64 bytes equal \texttt{hash\_only(content)} for the candidate content.
\end{itemize}

The \texttt{verify(raw, content)} function does exactly this verification: it parses the 73-byte raw form, extracts the trailing 64 bytes, recomputes \texttt{hash\_only(content)}, and returns a boolean. \texttt{verify} is the canonical way to ask ``does this \tesserahash{} reference this content?''

\subsection{The \texttt{hash\_at} function}

For deterministic test-vector generation and KAT reproducibility:

\begin{lstlisting}[language=Pseudocode]
raw = tessera_hash.hash_at(content, timestamp_ms=1700000000000, counter=0)
\end{lstlisting}

This bypasses the wall clock and the generator's counter state, producing a 73-byte \tesserahash{} whose header carries exactly the provided timestamp and counter. \texttt{hash\_at} MUST reject a \texttt{timestamp\_ms} outside $[0,\ 2^{48}-1]$ or a \texttt{counter} outside $[0,\ \code{0xFFFF}]$ with an error --- never by silently masking the value to the field width (Python raises \texttt{ValueError}; C returns \texttt{TESSERA\_HASH\_E\_RANGE}; Go returns \texttt{ErrTesseraHashRange}; the other languages raise their equivalent range error). Used by \texttt{tessera-hash/conformance/scripts/generate\_tessera\_hash.py} to produce the KAT files (\cref{sec:conformance}).

\subsection{What the provenance header tells you}

\begin{itemize}
\item \textbf{When} (timestamp): the millisecond a particular generator first produced this identifier. Two generators that independently hash the same content at different times will produce different timestamps but identical tails. The timestamp is a fact about the \emph{identifier}, not the content.
\item \textbf{In what order} (counter): the relative ordering of identifiers from a \emph{single} generator within a millisecond. Counters across generators are unrelated; a counter value of 7 from one generator means nothing about a counter value of 7 from another generator.
\item \textbf{Which content} (tail): the 512-bit SHAKE-256 hash of the input. This is the only field that carries content identity.
\end{itemize}

\textbf{What the provenance header does \emph{not} tell you:} who generated the identifier (no producer field; no signature), where it was generated (no node ID), or whether the generator's clock is honest (no clock attestation). Provenance is a self-claim of the generator; trust comes from the outer protocol.
